\section{Worked problems: a request is more than its text}
\begin{minipage}{\linewidth}
\subsection{One prefix, four outcomes}
\textbf{Problem.} A log contains only the generated bytes for a space followed
by \texttt{Rust}. Can you infer that the model selected EOS? What minimum
additional evidence distinguishes the four experiments?

\textbf{Solution.} No. EOS, the token budget, cancellation and a failure on
the next invocation can all follow that prefix. Record the request identity,
ordered event types and final outcome. To distinguish a selected EOS from a
budget stop, record the selection trace as well: the budget may prevent the
next model invocation entirely. Text alone is not a completion certificate.
\end{minipage}\par

\begin{minipage}{\linewidth}
\subsection{Absent is not zero}
\textbf{Problem.} The prompt \texttt{I} produces an immediate EOS and no text.
What should its first-text latency field contain?

\textbf{Solution.} No value. The event defining the endpoint never happened.
Zero would assert that text was emitted at arrival. Including zeros in an
average would make a workload with many empty completions appear faster.
Report the count of defined measurements and the count of empty completions.
\end{minipage}\par

\Needspace{12\baselineskip}
\begin{minipage}{\linewidth}
\subsection{Test a cancellation boundary}
\textbf{Problem.} Design a test for cancellation observed after a model call
returns but before its token is emitted.

\textbf{Solution rubric.} Use a controlled cancellation source and a model
whose invocation is observable. Assert one invocation, zero subsequent token
or text events, exactly one Cancelled terminal event and no later events.
The test must not rely on sleeping for a guessed duration. Such a test proves
a cooperative boundary, not interruption of a running kernel.
\end{minipage}\par

\Needspace{8\baselineskip}
\begin{minipage}{\linewidth}
\subsection{Price an abandoned cache}
\textbf{Problem.} Illustratively, each abandoned request holds a constant
576~MiB for 37~s. Abandonments arrive steadily at 0.5 per second.
Find the occupancy per request and mean resident memory. What fails if the
cache grows during decode?

\textbf{Solution.} The memory-time cost is
$576\times37=21{,}312$~MiB\,s, or 20.8125~GiB\,s per request.
There are $0.5\times37=18.5$ abandoned requests resident on average, costing
$18.5\times576/1024=10.40625$~GiB. With growing caches, replace the constant
product by the integral of resident bytes over the abandoned lifetime.
Do not multiply by a request fraction and call it a fleet utilization
fraction without matching duration and resource populations.
\end{minipage}\par

\Needspace{8\baselineskip}
\begin{minipage}{\linewidth}
\subsection{Why a full queue cannot carry its own rescue}
\textbf{Problem.} Data enters at 25 events/s and leaves at five, with
capacity 64. Estimate its fill time and payload bytes at 306 bytes/event.
Should cancellation wait to enqueue behind that data?

\textbf{Solution.} Net growth is 20 events/s, giving 3.2~s from empty.
The payload occupies $64\times306=19{,}584$ bytes, excluding representation
overheads. A cancellation sent through a full data queue can never arrive
if the consumer has gone. The teaching owner instead has an independent
stop operation and stores terminal metadata outside the data deque.
Real services should also bound serialized bytes, not only event count.
\end{minipage}\par

\Needspace{8\baselineskip}
\begin{minipage}{\linewidth}
\subsection{A terminal before retirement}
\textbf{Problem.} The owner submits ticket $(7,12,0)$, commits Cancelled
and delivers its terminal. The worker still holds the ticket. May slot seven
be reassigned? What should happen to its late reply?

\textbf{Solution.} Not yet. The terminal fixes output policy, not device
completion. A matching retirement clears the outstanding ticket and drops
the late payload. Only then is the modeled lease reusable. If generation
13 later receives a duplicate reply from generation 12, it must reject it
without clearing generation 13's ticket. A same-generation stale serial
must also fail. The lab tests both cases; a production device fence and
shared-block reference accounting remain separate obligations.
\end{minipage}\par

\Needspace{8\baselineskip}
\begin{minipage}{\linewidth}
\subsection{A wake is not our result}
\textbf{Problem.} An illustrative handler starts at 0~ms, waits in 100~ms
slices, and restarts after unrelated wakes every 4~ms through 2,000~ms.
When is its first timeout? What changes if its deadline was fixed at 100~ms?

\textbf{Solution.} The final wake establishes expiry at
$2{,}000+100=2{,}100$~ms. A client gone at 350~ms is therefore unobserved
for about 1,750~ms. A fixed initial deadline stays at 100~ms: at the 4~ms
wake, 96~ms remain. A periodic detector needs subsequent fixed checks to
notice the later disconnect; the first 100~ms check precedes it.
These are chosen integer times, not an operating-system timing guarantee.
\end{minipage}\par

\Needspace{8\baselineskip}
\begin{minipage}{\linewidth}
\subsection{Falsify a weak cancellation checker}
\textbf{Problem.} A three-token-budget request produces tokens 0, 1, 2 and
one MaxTokens terminal. The harness records a cancellation while the
request was still running after token 0. Why is checking the natural
outcome insufficient?

\textbf{Solution.} The counts match the uncancelled oracle but violate the
observed cancellation contract. In the step engine, cancellation precedes
the next numerical step, so the legal terminal is Cancelled after one token.
A checker must reject the natural trace, not accept it because it ends
once. This is the negative test
\texttt{checker\_rejects\_natural\_end\_after\_observed\_cancellation}.
Safety, natural stopping and cooperative cancellation are different checks.
\end{minipage}\par
